\setcounter{section}{19}

  \zwischenueberschrift{Soal latihan}

\inputaufgabe
{}
{

Tunjukkan bahwa
\definitionsverweis {matriks fundamental kedua}{}{}
bersifat
\definitionsverweis {simetris}{}{.}

}
{} {}

\inputaufgabe
{}
{

Misalkan
\maabbeledisp {\psi} {\R^2} {\R
} {(x,y)} { x^2-xy-y^3
} {,}
dan misalkan $Y$ adalah
\definitionsverweis {grafik}{}{}
dari $\psi$, yang dilengkapi dengan
\definitionsverweis {medan normal satuan}{}{}
berarah ke atas, serta
\maabbdisp {\varphi  =
\operatorname{Id}  \times \psi} {\R^2} {Y
} {}
adalah parametrisasi grafik tersebut. Tentukan
\definitionsverweis {matriks fundamental kedua}{}{}
untuk $\varphi$.

}
{} {}

\inputaufgabe
{}
{

Misalkan
\maabbeledisp {\psi} {\R^2} {\R
} {(x,y)} { e^{xy} -x^2y^3  \sin  \left(  x-y^2 \right)
} {,}
dan misalkan $Y$ adalah
\definitionsverweis {grafik}{}{}
dari $\psi$, yang dilengkapi dengan
\definitionsverweis {medan normal satuan}{}{}
berarah ke atas, serta
\maabbdisp {\varphi  =
\operatorname{Id}  \times \psi} {\R^2} {Y
} {}
adalah parametrisasi grafik tersebut. Tentukan
\definitionsverweis {matriks fundamental kedua}{}{}
untuk $\varphi$.

}
{} {}

\inputaufgabe
{}
{

Kita meninjau sfer satuan dengan
\definitionsverweis {medan normal satuan}{}{}
berarah ke luar dan parametrisasi
\maabbeledisp {\varphi} {[0, 2 \pi] \times [-1,1]} {\R^3
} {(u,v)} { \left(  \sqrt{1-v^2}  \cos  u   , \,   \sqrt{1-v^2}  \sin  u  , \,   v                 \right)
} {,}
lihat
[[Kugeloberfläche/Koordinaten von Zylinder aus/Horizontale Projektion/Flächenberechnung/Beispiel|Contoh 17.6]].
Tentukan
\definitionsverweis {matriks fundamental kedua}{}{}
untuk $\varphi$ pada interior domain parameter.

}
{} {}

\inputaufgabe
{}
{

Kita meninjau sfer satuan dengan
\definitionsverweis {medan normal satuan}{}{}
berarah ke luar dan parametrisasi
\maabbeledisp {\varphi} {[0, 2 \pi] \times \R} {\R^3
} {(u,v)} { { \frac{   1  }{  \sqrt{1+v^2}  } }  ( \cos  u  ,  \sin  u  , v )
} {,}
lihat
[[Kugeloberfläche/Koordinaten von Zylinder aus/Mittelpunktsprojektion/Flächenberechnung/Beispiel|Contoh 17.7]].
Tentukan
\definitionsverweis {matriks fundamental kedua}{}{}
untuk $\varphi$.

}
{} {}

\inputaufgabe
{}
{

Tentukan
\definitionsverweis {pemetaan Weingarten}{}{}
terhadap basis $\partial_1 \varphi, \partial_2 \varphi$ untuk berbagai parametrisasi $\varphi$ dari sfer satuan dalam
[[Kugeloberfläche/Koordinaten von Zylinder aus/Horizontale Projektion/Flächenberechnung/Beispiel|Contoh 17.6]],
[[Kugeloberfläche/Koordinaten von Zylinder aus/Mittelpunktsprojektion/Flächenberechnung/Beispiel|Contoh 17.7]],
dan
[[Kugeloberfläche/Geozentrische Koordinaten/Flächenberechnung/Beispiel|Contoh 17.8]].

}
{} {}

\inputaufgabe
{}
{

Tentukan
\definitionsverweis {kelengkungan Gauss}{}{}
sfer satuan dengan memakai berbagai parametrisasi $\varphi$ dari
[[Kugeloberfläche/Koordinaten von Zylinder aus/Horizontale Projektion/Flächenberechnung/Beispiel|Contoh 17.6]],
[[Kugeloberfläche/Koordinaten von Zylinder aus/Mittelpunktsprojektion/Flächenberechnung/Beispiel|Contoh 17.7]],
dan
[[Kugeloberfläche/Geozentrische Koordinaten/Flächenberechnung/Beispiel|Contoh 17.8]],
dengan menggunakan
[[Hyperfläche/Raum/Parametrisierung/Weingartenabbildung/Fundamentalformen/Fakt|Lema 19.3  (4)]].

}
{} {}

\inputaufgabe
{}
{

Misalkan $Y$ adalah suatu
\definitionsverweis {kurva}{}{}
di $\R^2$ yang dilengkapi dengan suatu medan normal satuan. Hubungkan
\definitionsverweis {kelengkungan}{}{}
kurva tersebut dengan
\definitionsverweis {matriks fundamental kedua}{}{.}

}
{} {}

\inputaufgabe
{}
{

Misalkan

\mavergleichskette
{\vergleichskette
{ Y
}
{ \subseteq }{ \R^3
}
{  }{
}
{    }{
}
{   }{
}
}
{}{}{}
adalah suatu
\definitionsverweis {permukaan berorientasi}{}{}
dan misalkan
\maabbdisp {\varphi} { V } { Y
} {,}

\mavergleichskette
{\vergleichskette
{ V
}
{ \subseteq }{ \R^2
}
{  }{
}
{    }{
}
{   }{
}
}
{}{}{,}
adalah parametrisasi lokal $Y$ yang terdiferensialkan tiga kali, dengan parameter $u,v$. Misalkan
\mathl{\left(  g_{ij}  \right)}{} adalah
\definitionsverweis {matriks fundamental pertama}{}{}
pada $V$. Tunjukkan bahwa
\definitionsverweis {kelengkungan Gauss}{}{}
memenuhi hubungan

\mavergleichskettedisp
{\vergleichskette
{ K
}
{ =} { \begin{aligned}[t]
& { \frac{   1  }{  \left(  g_{11} g_{22} -g_{12}^2  \right)^2 } }  \Biggl(
\\
& \det  \begin{pmatrix}  - { \frac{   1  }{  2 } } \partial_2^2  g_{11}  + \partial_1 \partial_2 g_{12} - { \frac{   1  }{  2 } }  \partial_1^2  g_{22}   &   { \frac{   1  }{  2 } } \partial_1 g_{11}  &  \partial_1 g_{12} - { \frac{   1  }{  2 } } \partial_2 g_{11}  \\  \partial_2 g_{12} -{ \frac{   1  }{  2 } } \partial_1 g_{22} &  g_{11}  &  g_{12}  \\ { \frac{   1  }{  2 } } \partial_2 g_{22}  &  g_{12}  &  g_{22}  \end{pmatrix}
\\
& {}- \det  \begin{pmatrix}  0   &   { \frac{   1  }{  2 } }  \partial_2 g_{11}  &  { \frac{   1  }{  2 } }  \partial_1 g_{22}   \\  { \frac{   1  }{  2 } } \partial_2 g_{11} & g_{11} & g_{12} \\ { \frac{   1  }{  2 } } \partial_1 g_{22}  &  g_{12} & g_{22} \end{pmatrix} \Biggr)
\end{aligned}
}
{ } {
}
{ } {
}
{ } {
}
}
{}{}{}
tersebut.

}
{} {}

  \zwischenueberschrift{Soal untuk dikumpulkan}

\inputaufgabe
{3}
{

Misalkan
\maabbeledisp {\psi} {\R^2} {\R
} {(x,y)} { x^3-2x^2y-7xy^3+x^2y^2+5y^4
} {,}
dan misalkan $Y$ adalah
\definitionsverweis {grafik}{}{}
dari $\psi$, yang dilengkapi dengan
\definitionsverweis {medan normal satuan}{}{}
berarah ke atas, serta
\maabbdisp {\varphi  =
\operatorname{Id}  \times \psi} {\R^2} {Y
} {}
adalah parametrisasi grafik tersebut. Tentukan
\definitionsverweis {matriks fundamental kedua}{}{}
untuk $\varphi$.

}
{} {}

\inputaufgabe
{4}
{

Kita meninjau sfer satuan dengan
\definitionsverweis {medan normal satuan}{}{}
berarah ke luar dan parametrisasi
\maabbeledisp {\varphi} {{]{- { \frac{   \pi }{  2 } }},  { \frac{   \pi }{  2 } } [} \times {]{-\pi}, \pi[} } {\R^3
} {(u,v)} {(\cos  u  \cos  v  , \cos  u  \sin  v  , \sin  u )
} {,}
lihat
[[Kugeloberfläche/Geozentrische Koordinaten/Flächenberechnung/Beispiel|Contoh 17.8]].
Tentukan
\definitionsverweis {matriks fundamental kedua}{}{}
untuk $\varphi$.

}
{} {}

\inputaufgabe
{5 (2+1+1+1)}
{

Misalkan
\maabb {\gamma} {]a,b[} { \R^2
} {}
adalah suatu kurva injektif yang
\definitionsverweis {terdiferensialkan dua kali}{}{}
dan
\definitionsverweis {berparameter panjang busur}{}{,}
dengan kurva citra

\mavergleichskettedisp
{\vergleichskette
{ C
}
{ \subseteq} { V
}
{ \subseteq} { \R^2
}
{ } {
}
{ } {
}
}
{}{}{,}
yang tertutup dalam himpunan terbuka $V$. Kita meninjau silinder di atas kurva ini, yakni pemetaan
\maabbdisp {\varphi = \gamma \times
\operatorname{Id}_{ \R }} {]a,b[ \times \R } { C \times \R
} {.}
Lengkapi silinder tersebut dengan salah satu dari kedua medan normal satuannya.
\aufzaehlungvier{Hitung
\definitionsverweis {matriks fundamental pertama}{}{}
dan
\definitionsverweis {matriks fundamental kedua}{}{}
dari $\varphi$.
}{Hitung
\definitionsverweis {pemetaan Weingarten}{}{}
terhadap basis $\partial_1 \varphi,  \partial_2 \varphi$.
}{Hitung
\definitionsverweis {kelengkungan-kelengkungan utama}{}{}
dari $C \times \R$.
}{Hitung
\definitionsverweis {kelengkungan Gauss}{}{}
dari $C \times \R$.
}

}
{} {}

 [[Kategorie:Latexseite]]
